% Figure: constrained maximum and minimum of a linear objective on an
% ellipse. Maximise and minimise f(x,y) = x + 2y on the ellipse
% x^2/4 + y^2 = 1. The two optima are the points where a level line of f
% (slope -1/2) is tangent to the ellipse; the far tangent is the maximum,
% the near tangent the minimum, and the two constraint gradients point in
% opposite senses, so the multiplier changes sign between them.
\begin{figure}[htbp]
\centering
\begin{tikzpicture}[scale=1.0,
  ell/.style={engblue, very thick},
  lev/.style={enggray, thin},
  tang/.style={engred, thick},
  gvec/.style={engorange, very thick, -{Stealth[length=2.6mm]}},
  pt/.style={circle, fill=engblack, inner sep=1.5pt},
]
  \begin{axis}[
    width=9.6cm, height=6.8cm,
    axis equal image,
    xlabel={$x$}, ylabel={$y$},
    xmin=-3.4, xmax=3.6, ymin=-1.9, ymax=2.0,
    axis lines=middle,
    xtick={-3,-2,-1,1,2,3}, ytick={-1,1},
    tick label style={font=\small},
    label style={font=\small},
    clip=false,
  ]
    % The ellipse x^2/4 + y^2 = 1.
    \addplot[ell, domain=0:360, samples=100] ({2*cos(x)}, {sin(x)});
    % A few level lines f = x + 2y = c, slope -1/2. Non-tangent guides.
    \foreach \c in {-2,0} {
      \addplot[lev, domain=-3.4:3.6, samples=2] {(\c - x)/2};
    }
    % Optimum: grad f = (1,2). f max = sqrt(1+16)= sqrt(17) approx 4.123 at
    % (x,y) = (2*1/sqrt(17*... )). Solve: on ellipse, extremum of x+2y.
    % param x=2cos t, y=sin t; f = 2cos t + 2 sin t = 2 sqrt2 sin(t+45).
    % max at t=45deg: x=2/sqrt2=1.414, y=1/sqrt2=0.707, f=2 sqrt2=2.828.
    % min at t=225deg: x=-1.414, y=-0.707, f=-2.828.
    % Tangent lines through those points, slope -1/2.
    \addplot[tang, domain=-1.2:3.6, samples=2] {(2.828 - x)/2};
    \addplot[tang, domain=-3.4:1.2, samples=2] {(-2.828 - x)/2};
    \node[pt] at (axis cs:1.414,0.707) {};
    \node[pt] at (axis cs:-1.414,-0.707) {};
    \node[engred, anchor=south west, font=\small] at (axis cs:1.5,0.78)
      {max};
    \node[engred, anchor=north east, font=\small] at (axis cs:-1.5,-0.78)
      {min};
    % Constraint gradient (outward normal) at the max point ~ direction (x/2, 2y).
    \draw[gvec] (axis cs:1.414,0.707) -- (axis cs:2.05,2.12);
    \node[engorange, anchor=west, font=\small] at (axis cs:2.0,1.6)
      {$\nabla g$};
  \end{axis}
\end{tikzpicture}
\caption[Constrained maximum and minimum of a linear objective on an
ellipse.]{The objective $f(x, y) = x + 2y$ has straight level lines of slope
$-\tfrac{1}{2}$ (grey). On the ellipse $x^{2}/4 + y^{2} = 1$ (blue) the
objective is extremised where a level line runs tangent to the curve (red).
The far tangent gives the constrained maximum $f = 2\sqrt{2}$ at
$(\sqrt{2}, \tfrac{1}{\sqrt{2}})$; the near tangent gives the constrained
minimum $f = -2\sqrt{2}$ at the antipodal point. At both, the objective
gradient $(1, 2)$ is parallel to the constraint gradient (orange), which is
the Lagrange condition; the multiplier is positive at the maximum and
negative at the minimum, because the two constraint normals point in
opposite senses relative to $\grad f$. One constraint yields two optima, the
generic feature of extremising over a closed curve.}
\label{fig:vol02:ch11:ellipse-extrema}
\end{figure}
